Section 1.5 has partial formalization in its adjacent Lean companion and the reusable module \texttt{TraceFormulas.Geometry.IsometryTypes}. The manuscript first classifies fixed points, then gives separate rotation, translation, and dilation normal forms, and finally treats the three translation lengths separately. The hyperbolic diagram illustrates the semicircular axis and its conjugate imaginary axis, with corresponding points related by $h(g\cdot z)=e^r h(z)$. Its arclength interpretation uses the existing vertical distance formula. The generator checks its plotted coordinates; this is separate from the outstanding Lean proofs of geometric conjugacy, path distance, and intrinsic axis characterization.

For the \hyperref[hyperbolic.types]{type discriminant}, Lean uses Mathlib's actual matrix trace and proves invariance under replacing $g$ by $-g$. For the \hyperref[hyperbolic.fixed-point-classification]{fixed-point classification}, it proves the determinant-one discriminant calculation and the exact fixed-point equation \eqref{eq:mobius-fixed-point} for points of Mathlib's upper half-plane. It also derives $D(g)=-4(c\operatorname{Im}z)^2$ from an interior fixed point. This last identity is an additional algebraic ingredient, not the full classification: it does not supply a fixed point from the sign of $D(g)$ or count projective boundary fixed points.

For the \hyperref[hyperbolic.elliptic-normal-form]{elliptic normal form}, the companion proves that an already supplied fixed point $p$ gives $s(p)^{-1}gs(p)=k(\theta)$ for some $\theta$. It follows the manuscript's conjugation and stabilizer argument. Existence and uniqueness of $p$ from the elliptic hypothesis, exclusion of central rotations, and the Cayley rotation identity remain to be connected in Lean. The \hyperref[hyperbolic.parabolic-normal-form]{parabolic} and \hyperref[hyperbolic.hyperbolic-normal-form]{hyperbolic normal-form propositions} are not yet formalized. The checked statements show that $n(t)$ fixes an interior point exactly when $t=0$, and $a(y)$ fixes one exactly when $y=1$; these include the identity cases without claiming that they are parabolic or hyperbolic.

The calculation in the \hyperref[hyperbolic.parabolic-length]{parabolic translation-length proof} is checked as $u(z,n(t)\cdot z)=t^2/(4(\operatorname{Im}z)^2)$. For the \hyperref[hyperbolic.hyperbolic-length]{hyperbolic translation-length proof}, Lean proves
\[
 u(z,a(y)\cdot z)
 =\frac{(y-1)^2}{4y}\left(1+\frac{(\operatorname{Re}z)^2}{(\operatorname{Im}z)^2}\right),
 \qquad y>0.
\]
It proves the lower bound $(y-1)^2/(4y)$ and, for $y\ne1$, equality exactly on the imaginary axis. Here $y$ is the positive dilation parameter, not the imaginary part of $z$; the manuscript substitutes $y=e^r$. This modest generalization also covers contracting dilations and keeps the identity case explicit.

The \hyperref[hyperbolic.translation-length]{metric-infimum definition}, its \hyperref[hyperbolic.translation-length-conjugacy]{conjugacy invariance}, all three translation-length conclusions, attainment or nonattainment, the trace formula for the length, and the intrinsic geodesic-axis characterization remain open. In particular, minimization of $u$ is not presented as a proof of minimization of the prose's path distance: the bridge between the established coordinate invariant and that distance still has to be formalized. No infimum or trace formula is defined to equal its desired answer.

The machine-check snapshot is recorded in \texttt{lean/audits/chapter01-release-check-2026-11-12/record.json}. Its scope is partial Chapter 1 coverage. Steven's review of mathematical correspondence remains pending independently of the Lean build and axiom audit.
